\section{Discussion}
\label{sec:discussion}

\textbf{Breaking the Limits of Linear Spatial Filtering.} 
The classification of intra-limb movements presents a severe challenge due to the heavily overlapping cortical activation patterns within the same hemisphere. As demonstrated in Table~\ref{tab:main}, the classical exploratory CSP pipeline hovered near chance-level accuracy ($\sim$36\%) for this four-class task %[cite: 79]. 
The transition to convolutional receptive fields (Block II) yielded an immediate improvement of approximately 20 percentage points over the optimized shallow CSP baseline %[cite: 6]. 
This substantial leap suggests that successful same-limb decoding relies heavily on non-linear spatio-temporal dynamics that simple spatial variance filters inherently fail to capture %[cite: 7]. 
However, it is worth noting that for complete parity in future benchmarking, CSP and CNN pipelines must strictly unify band-pass envelopes, referencing, and temporal cropping before blaming spatial filtering deficits alone %[cite: 10].

\textbf{The Efficacy and Boundaries of Multimodal Fusion.}
While deep temporal filters began untangling the latent manifolds, synergistic movements (e.g., hand versus forearm) still suffered from significant misclassification bottlenecks %[cite: 71]. 
Integrating multimodal priors successfully addressed this. Welch PSD concatenation and cross-spectral coherence attention each added roughly 5 percentage points of accuracy over the EEGNet-style baseline %[cite: 11]. 
The $t$-SNE embeddings in Figure~\ref{fig:master_grid} visually confirm that fusion forces tighter intra-class clustering and resolves boundary ambiguities %[cite: 72, 73]. 
Nevertheless, because withheld softmax inputs from different classes remain neighboring points across all Block II models %[cite: 15], 
we conclude that auxiliary maps supply marginal, complementary constraints rather than entirely orthogonal, linearly separable geometry %[cite: 12].

\textbf{Neurophysiological Validity versus Raw Accuracy.}
In neurotechnology, interpretable connectivity structures and final softmax margins are not interchangeable metrics %[cite: 14]. 
While Table~\ref{tab:main} shows only modest incremental accuracy gains from connectivity fusion %[cite: 13], 
Figure~\ref{fig:connTopo} proves that the network's attention mechanisms are physically grounded. The organized peri-central contrast between trial-averaged hand epochs and the resting state confirms that the fusion network relies on genuine sensorimotor structures rather than overfitting to hardware noise or muscular artifacts %[cite: 13, 88].

\textbf{Generalization and Future Work.}
While this study establishes the superiority of deep multimodal fusion for complex upper-limb decoding, the current validation strategy relies on withholding trials from a pooled subject mix %[cite: 16]. 
This approach is inherently optimistic relative to the strict leave-one-subject-out (LOSO) cross-validation required for clinical neurotechnology translation %[cite: 16]. 
Recommended follow-ups include strictly subject-independent benchmarks to evaluate true model generalization %[cite: 4], 
alongside the exploration of fusion designs that exploit auxiliary tensors without weakening the underlying temporal discriminability %[cite: 4].